\scanpage{25}
d'écrire que $u_K$ commute \unsure{:|à} $\mathfrak{g}_K$, or par hyp. il commute aux $\omega(X)\in\mathfrak{g}_K$ $(X\in\operatorname{Dér}(\unsure{H|\Pi|K}/k_0)\unsure{]|)}$, Soit $\mathfrak{u}\subset\mathfrak{g}$ le plus petit sous-espace vectoriel de $\mathfrak{g}$ tel que $\mathfrak{u}_K$ contienne les $\omega(X)$. Alors il est clair (alg. linéaire) que \deleted{$\mathfrak{u}$} $u$ commute \unsure{:|à} $\mathfrak{u}$, donc il suffit que $\mathfrak{u}=\mathfrak{g}$ pour pouvoir conclure (N.B. Il suffirait que l'\unsure{algèbre|Algèbre|algébre} de Lie engendrée par $\mathfrak{u}$ dans $\mathfrak{g}$ soit $\mathfrak{g}$, mais en fait on se convainc facilement que $\mathfrak{u}$ est déjà une $\unsure{k|K}$-algèbre de Lie....). On peut $\overline{\mathrm m}$ remplacer $\mathfrak{u}$ par l'algèbre de Lie algébrique qu'elle engendre
% archive sheet number: 24
